% УВОД, без номер. Проблемът и цената му вървят преди архитектурата.
\section*{УВОД}
\label{sec:introduction}
\addcontentsline{toc}{section}{УВОД}

Всяка организация с информационни технологии има повече идеи за проекти,
отколкото пари и хора да ги изпълни. Кой се финансира и в какъв ред решава
управлението. На практика решението се взема в електронна таблица, а таблицата
има три недостатъка, които не са естетически.

Първо, допусканията в нея не са записани никъде. Кой критерий колко тежи, защо
един проект е предпоставка за друг, откъде идва разчетната стойност: всичко живее
в главата на попълнилия клетките, и шест месеца по-късно никой, включително той,
не може да каже защо е така.

Второ, отхвърлянето идва без причина. Не се казва дали проектът е спрян от
бюджета, от липсата на инженери, от неизпълнена предпоставка или от по-ниската
оценка. Разликата е съществена: бюджетът се преразглежда идната година, а ниската
оценка не се.

Трето, и това е най-скъпото, таблицата смята с очаквани стойности, при които
портфейл, побиращ се в бюджета, изглежда безопасен. Настоящата разработка измери
какъв е той в действителност върху съставен сценарий: портфейл с очаквана стойност
1123 при бюджет 1150 има \textbf{38,7 процента} вероятност да остане в бюджета,
когато входните оценки се изтеглят от собствените си разпределения. Тоест около
три на пет шанс за преразход при решение, което изглежда с резерв. Таблицата с
очаквани стойности не може да покаже това число.

Проектът разработва система, наречена \textbf{Quorum}, която подпомага
решението, вместо да го замества. Тя приема бизнес процеса като насочен граф от
дейности със стойности, продължителности и зависимости, избира портфейла с
най-висока сумарна полза при зададените ограничения и преизчислява избора при
изменени допускания. Трите недостатъка определят трите изисквания: допусканията са
в текстов файл, който се версионира; всеки отхвърлен проект се връща с
ограничението, което го е спряло; рискът се смята върху разпределения.

\textbf{Целта} е класически методи на управленската наука, многокритериален анализ
и целочислена оптимизация, да бъдат вградени в работеща система така, че
препоръката да е проследима до входните допускания. \textbf{Обхватът} спира при
одобряването на портфейла: не са реализирани проследяване на изпълнението,
счетоводно отчитане, интеграция с ERP и уеб интерфейс.
